\documentclass[portrait,a0paper]{baposter}

\usepackage{calc}
\usepackage{amsmath}
\usepackage{amssymb}
\usepackage{relsize}
\usepackage{multirow}
\usepackage{rotating}
\usepackage{bm}
\usepackage{url}
\usepackage{graphicx}
\usepackage{multicol}
\usepackage{soul} % for strikethrough
\usepackage{color}

% Custom page size is handled by baposter class options
% \usepackage[paperwidth=80cm,paperheight=100cm,margin=2cm]{geometry}

\definecolor{darkblue}{rgb}{0.0,0.08,0.45}
\definecolor{lightblue}{rgb}{0.8,0.85,0.95}

\begin{document}

\begin{poster}
{
  grid=false,
  headerborder=closed,
  colspacing=1em,
  columns=2,
  bgColorOne=white,
  bgColorTwo=white,
  borderColor=darkblue,
  headerColorOne=lightblue,
  headerColorTwo=lightblue,
  headerFontColor=darkblue,
  boxColorOne=white,
  textborder=rounded,
  eyecatcher=true,
  headershape=rounded,
  headershade=shadelr,
  headerfont=\Large\bf\textsc,
  textfont={\setlength{\parindent}{0em}\large},
  boxshade=plain,
  background=plain,
  linewidth=2pt
}
% Eye Catcher (left of title)
{}
% Title
{\bf\textsc{\st{How} Can We Infer Contact Matrices?}\vspace{0.3em}}
% Authors
{\textsc{Michael Klots, Yair Daon$^\star$}\\[0.3em]
  \textit{Azrieli Faculty of Medicine, Bar-Ilan University, Safed, Israel}\\[0.5em]
$^\star$yair.daon@gmail.com}
{}

\headerbox{Background}{name=map,column=0,row=0,span=2}{
  When studying spatial spread of infectious diseases we need data
  on connectivity e.g.~flights, commutes (Daon et al.
  2020).

  \textbf{Problem: Instead of relying on connectivity data, we want to infer it from seasonal outbreak patterns, e.g.~flu.}
  
  \begin{center}
    \includegraphics[width=\linewidth,height=0.15\textheight]{../pix/flights.pdf}
                    {\small Data: (Mao et al., 2015)}
  \end{center}
}


%%%%%%%%%%%%%%%%%%%%%%%%%%%%%%%%%%%%%%%%%%%%%%%%%%%%%%%%%%%%%%%%%%%%%%%%%%%%%%
\headerbox{Seasonally forced multi-region model}{name=model,column=0,below=map}{
%%%%%%%%%%%%%%%%%%%%%%%%%%%%%%%%%%%%%%%%%%%%%%%%%%%%%%%%%%%%%%%%%%%%%%%%%%%%%%
\vspace{-0.5em}
  \begin{align*}
    \dot{\mathbf{S}} &= \mu - {\beta(t) \, \mathbf{S} \circ \mathbf{C} \mathbf{I}} - \mu \mathbf{S} \\
    \dot{\mathbf{I}} &= {\beta(t) \, \mathbf{S} \circ \mathbf{C} \mathbf{I}} - \gamma \mathbf{I} - \mu \mathbf{I} \\
    %% \dot{\mathbf{I}} &= \sigma \mathbf{E} - \gamma \mathbf{I} - \mu \mathbf{I} \\
    \dot{\mathbf{R}} &= \gamma \mathbf{I} - \mu \mathbf{R}
  \end{align*}
  Contact matrix $\mathbf{C}$ has entries in $[0,1]$.%; seasonal driver $\beta(t)$.
%% \vspace{0.1em}
}

%%%%%%%%%%%%%%%%%%%%%%%%%%%%%%%%%%%%%%%%%%%%%%%%%%%%%%%%%%%%%%%%%%%%%%%%%%%%%%
\headerbox{Inference for contact matrix}{name=inverse, column=0, below=model}{
%%%%%%%%%%%%%%%%%%%%%%%%%%%%%%%%%%%%%%%%%%%%%%%%%%%%%%%%%%%%%%%%%%%%%%%%%%%%%%
  \vspace{-0.5em}
  %Incidence satisfies:%Assume $S(t_i), I(t_i)$ for $t_i, i=1,\ldots,T$ known (unrealistic). 
  \begin{align*}
    \text{incidence}_i \approx \beta(t_i) \, \mathbf{S}(t_i) \circ \mathbf{C} \mathbf{I}(t_i)
    \Longrightarrow \mathbf{C} \mathbf{I}(t_i) \approx \frac{\text{incidence}_i}{\beta(t_i) \mathbf{S}(t_i)}
  \end{align*}
  
  %Define:%Collect $\mathbf{I}(t_i)$ into a matrix:
  \begin{equation*}
    M := \left[
      \begin{array}{cccc}
        \vert & \vert &        & \vert \\
        \mathbf{I}(t_1)    & \mathbf{I}(t_2)    & \ldots & \mathbf{I}(t_T)    \\
        \vert & \vert &        & \vert 
      \end{array}
      \right]
  \end{equation*}
  
  %\vspace{1em}
  \begin{equation*}
    {\Large\textbf{Then:}\ } \boxed{\color{red}\text{relative error}(\textbf{C}) \leq \text{relative error}(\text{incidence}) \kappa(M)}
  \end{equation*}
  
  $\kappa(M)$ condition number: ratio of max \& min singular values.
  %%\vspace{0.5em}
}


%%%%%%%%%%%%%%%%%%%%%%%%%%%%%%%%%%%%%%%%%%%%%%%%%%%%%%%%%%%%%%%%%%%%%%%%%%%%%%
\headerbox{Inference Fails. Why?}{name=results,column=1, below=map}{
%%%%%%%%%%%%%%%%%%%%%%%%%%%%%%%%%%%%%%%%%%%%%%%%%%%%%%%%%%%%%%%%%%%%%%%%%%%%%%
  Condition numbers:
  \begin{itemize}
  \item 2 regions: $\kappa(M) \approx {\color{red}10^5}$.
  \item 5 regions: $\kappa(M) \approx {\color{red}10^8}$.
  \item 10 regions: $\kappa(M) \approx {\color{red}10^{11}}$.
  \end{itemize}
  For two regions with 1\% relative error in incidence:
  \begin{equation*}
    \text{relative error}(\textbf{C}) \approx 0.01 \cdot 10^5 = 1000.
  \end{equation*}
  Solving the inverse problem for $\mathbf{C}$ (lower box) fails miserably. What is going on?
\begin{itemize}
\item seasonal forcing $\Rightarrow$ synchronized epidemics $\Rightarrow$ huge $\kappa(M)$.
\item \textbf{Inferring contact matrices from synchronized epidemic data is extremely ill-conditioned}.
\item Synchrony between regions makes the inverse problem intractable.
\item Traditional connectivity data remains essential.
\end{itemize}
%\vspace{0.25em}
}


\headerbox{Optimization Paths}{name=res2, below=inverse, span=2}{
  %%
  We try to solve the inverse problem for $\textbf{C}$ with standard
  methods. When seasonal driver is present (left) optimization does not
  converge. Without seasonal driver (right) true connectivity is
  found.
  \begin{center}
    \includegraphics[width=0.75\linewidth, height=0.36\linewidth]{../pix/poster.jpg}
  \end{center}
  \vspace{-0.1em}
}

%%%%%%%%%%%%%%%%%%%%%%%%%%%%%%%%%%%%%%%%%%%%%%%%%%%%%%%%%%%%%%%%%%%%%%%%%%%%%%
%% \headerbox{References}{name=refs,column=1,below=conclusion,span=1}{
%% %%%%%%%%%%%%%%%%%%%%%%%%%%%%%%%%%%%%%%%%%%%%%%%%%%%%%%%%%%%%%%%%%%%%%%%%%%%%%%
%% \small
%% \begin{thebibliography}{9}
%% \bibitem{mao2015} Mao et al. \textit{Journal of Transport Geography} (2015)
%% \bibitem{daon2020} Daon et al. \textit{Journal of Travel Medicine} (2020)  
%% \bibitem{trefethen} Trefethen \& Bau. Numerical Linear Algebra (1997)
%% \end{thebibliography}
%% }


\end{poster}

\end{document}
